\section{A planar corollary}
\label{sec:plane}

The full-Fock-space estimate also gives a gap on the plane. We use the
analytic lowest-Landau-level representation described by Girvin and
Jach~\cite[Section~II]{GirvinJach}, with Gaussian scale fixed below.
Let
$\mathcal B$ be the Bargmann space of entire functions with norm
\[
 \|f\|_{\mathcal B}^2=\frac1\pi\int_{\mathbb C}|f(z)|^2e^{-|z|^2}\,d^2z,
 \qquad e_j(z)=\frac{z^j}{\sqrt{j!}}.
\]
Multiplication by $\pi^{-1/2}e^{-|z|^2/2}$ identifies this space unitarily
with the planar lowest Landau level at magnetic field $2$; changing magnetic
length is a unitary dilation.
For a pair of coordinates set $w=(z_1-z_2)/\sqrt2$ and
$Z=(z_1+z_2)/\sqrt2$.  Let $P^{(1)}$ be the orthogonal projector onto
relative degree one, with no restriction on the degree in $Z$, and put
\[
 H_N^{\mathrm{pl}}=\left.\sum_{1\le i<j\le N}P_{ij}^{(1)}
                         \right|_{\bigwedge^N\mathcal B}.
\]
Thus each pair projector has coefficient one, as on the sphere.

\begin{corollary}[Planar spectral gap]\label{cor:plane}
For every $N\ge2$, the untruncated planar Hamiltonian satisfies
\[
 (H_N^{\mathrm{pl}})^2\ge\gamma_* H_N^{\mathrm{pl}},
\]
where $\gamma_*$ is the constant in Theorem~\ref{thm:fock}.
Equivalently, if $P_{0,N}^{\mathrm{pl}}$ projects onto its kernel, then
$H_N^{\mathrm{pl}}\ge\gamma_*(I-P_{0,N}^{\mathrm{pl}})$.
The constant is independent of both particle number and total angular
momentum.  In particular, the case $\ell=3$ of the spectral gap conjecture
in \cite[Conjecture~A.1]{Rougerie} holds, in its pair-projector normalization.
\end{corollary}
\begin{proof}
The coordinate change $(z_1,z_2)\mapsto(w,Z)$ is unitary for the
Bargmann inner product.  Consequently the vectors
$wZ^p/\sqrt{p!}$, $p\ge0$, form an orthonormal basis of the relative-degree-one
subspace.  Their coefficients in the normalized wedge basis are
\[
 v_p^{\mathrm{pl}}(i,j)
   =(i-j)\sqrt{\frac{p!}{2^p i!j!}}\,\mathbf1_{i+j=p+1},
 \qquad i<j.
\]
Indeed, expanding $(z_1-z_2)(z_1+z_2)^p$
first gives the coefficient in the tensor basis; passage to the normalized
antisymmetric wedge multiplies it by $\sqrt2$.  These are exactly the
limits in \eqref{eq:pairstar}.  The second-quantized pair-projector identity
therefore gives
$H_N^{\mathrm{pl}}=\sum_{p\ge0}B(v_p^{\mathrm{pl}})^\dagger B(v_p^{\mathrm{pl}})$
on antisymmetric polynomials.

Fix $N$ and a total degree $L$.  The corresponding homogeneous subspace
has the finite orthonormal basis
\[
 e_{i_1}\wedge\cdots\wedge e_{i_N},\qquad
 0\le i_1<\cdots<i_N,\qquad \sum_{k=1}^N i_k=L.
\]
For $Q\ge\max\{L,N-1,2\}$ this is also the entire spherical subspace of total
lowering deficit $L$, under the identification of normalized one-particle
bases.  Both Hamiltonians preserve this subspace: a pair is annihilated
and created with the same sum of orbital labels.  Only $p+1\le L$
can contribute.  The finitely many matrix entries of the spherical
restriction therefore converge to those of the planar restriction by
\eqref{eq:pair}--\eqref{eq:pairstar}; the convergence is in operator norm.

Fix $0<\gamma<\gamma_*$.  Theorem~\ref{thm:fock} gives
$H_Q^2\ge\gamma H_Q$ for all sufficiently large $Q$.
Restrict to the invariant block just described and let $Q\to\infty$.
Norm convergence passes through the square and gives
$(H_N^{\mathrm{pl}})^2\ge\gamma H_N^{\mathrm{pl}}$ on that block.
Now let $\gamma\uparrow\gamma_*$; the cone of positive matrices is closed,
so the same inequality holds at $\gamma_*$.  The limit $Q\to\infty$ is taken
separately for each fixed $\gamma$, allowing the flux threshold to depend
on $\gamma$.  Here $N$ and $L$ were arbitrary, and the resulting constant is the
same for all of them.

Finally, $\bigwedge^N\mathcal B$ is the Hilbert direct sum of its homogeneous
subspaces, and $0\le H_N^{\mathrm{pl}}\le\binom N2 I$.  The block inequalities
thus extend by continuity to the whole space.  The spectral theorem gives
the asserted gap.  Restriction to total degrees
$L\le3N(N-1)/2$ gives exactly the sector in the cited conjecture for
$\ell=3$.
\end{proof}
